% ============================================================
% Graphiques d'évaluation du système RAG
% À inclure dans un chapitre avec : % ============================================================
% Graphiques d'évaluation du système RAG
% À inclure dans un chapitre avec : % ============================================================
% Graphiques d'évaluation du système RAG
% À inclure dans un chapitre avec : % ============================================================
% Graphiques d'évaluation du système RAG
% À inclure dans un chapitre avec : \input{graphes-evaluation}
% ============================================================

\begin{figure}[htbp]
\centering
\begin{tikzpicture}[x=10cm,y=0.62cm]
  \draw[->,muted] (-0.1,0) -- (0.62,0)
    node[below=4pt,font=\scriptsize] {Score de pertinence maximal};
  \foreach \x/\label in {-0.1/{-0,10},0/{0},0.2/{0,20},0.4/{0,40},0.6/{0,60}} {
    \draw[muted] (\x,0.08) -- (\x,-0.08);
    \node[below=9pt,font=\scriptsize] at (\x,0) {\label};
  }
  \draw[accent,fill=accent!75] (0,0.45) rectangle (0.082,0.80);
  \draw[accent,fill=accent!75] (0,1.45) rectangle (0.115,1.80);
  \draw[accent,fill=accent!75] (0,2.45) rectangle (0.194,2.80);
  \draw[accent,fill=accent!75] (0,3.45) rectangle (0.567,3.80);
  \draw[warn,fill=warn!75] (0,4.45) rectangle (0.046,4.80);
  \draw[warn,fill=warn!75] (0,5.45) rectangle (0.003,5.80);
  \draw[warn,fill=warn!75] (-0.048,6.45) rectangle (0,6.80);
  \foreach \y/\label in {
    0.625/{Commande 10248},
    1.625/{Facture 10248},
    2.625/{Stock juillet},
    3.625/{Maria Larsson},
    4.625/{Congés payés},
    5.625/{Note de frais},
    6.625/{Intégration}
  } {
    \node[anchor=east,font=\scriptsize] at (-0.05,\y) {\label};
  }
  \node[anchor=west,font=\scriptsize] at (0.59,0.625) {0,082};
  \node[anchor=west,font=\scriptsize] at (0.59,1.625) {0,115};
  \node[anchor=west,font=\scriptsize] at (0.59,2.625) {0,194};
  \node[anchor=west,font=\scriptsize] at (0.59,3.625) {0,567};
  \node[anchor=west,font=\scriptsize] at (0.59,4.625) {0,046};
  \node[anchor=west,font=\scriptsize] at (0.59,5.625) {0,003};
  \node[anchor=west,font=\scriptsize] at (0.59,6.625) {-0,048};
\end{tikzpicture}
\caption{Scores maximaux de pertinence selon le type de question}
\label{fig:scores-retrieval}
\end{figure}

\begin{table}[htbp]
\centering
\caption{Évaluation récapitulative du système RAG}
\label{tab:evaluation-recapitulative}
\small
\begin{tabularx}{\textwidth}{@{}p{3.2cm}p{4.1cm}X p{2.1cm}@{}}
\toprule
\textbf{Dimension} & \textbf{Indicateurs} & \textbf{Résultats observés} &
\textbf{Appréciation} \\
\midrule
\shortstack[l]{\textbf{Qualité de la recherche}\\\eng{Retrieval}}
&
Score de pertinence, recherche top-$k$, couverture des sources
&
Recherche vectorielle fonctionnelle avec $k=4$ et affichage des scores.
Score moyen français : 0,153 ; score moyen anglais : 0,461. La recherche
par identifiant exact reste perfectible.
&
\textbf{Satisfaisante, à améliorer}
\\
\midrule
\shortstack[l]{\textbf{Qualité des réponses}\\\eng{Generation}}
&
Fidélité au contexte, pertinence, citations, maîtrise des hallucinations
&
Prompt restrictif, température fixée à 0,1, réponse de repli si le contexte
est absent et citation des fichiers sources. Les métriques RAGAS complètes
restent à mesurer.
&
\textbf{Satisfaisante}
\\
\midrule
\shortstack[l]{\textbf{Performance}\\\textbf{technique}}
&
Latence de recherche, latence de génération, capacité d'indexation
&
2\,676 documents traités et 2\,929 segments indexés. Latence de recherche :
182 ms en moyenne. Génération locale : environ 110 s sur CPU, ce qui
constitue le principal goulot d'étranglement.
&
\textbf{Satisfaisante}
\\
\midrule
\shortstack[l]{\textbf{Fiabilité et}\\\textbf{robustesse}}
&
Extraction, traçabilité, gestion des erreurs, reproductibilité
&
0 échec d'extraction observé sur le corpus mesuré. Métadonnées conservées,
sources et scores affichés, API structurée avec endpoints de santé,
d'interrogation et d'ingestion. Authentification et re-ranking à ajouter
avant une exposition publique.
&
\textbf{Bonne en local}
\\
\midrule
\shortstack[l]{\textbf{Utilité}\\\textbf{métier}}
&
Accès à l'information, confidentialité, simplicité d'utilisation,
vérifiabilité
&
Interface conversationnelle accessible en langage naturel, réponses
accompagnées de sources vérifiables et traitement local des documents.
L'outil est particulièrement adapté à une démonstration et à une recherche
documentaire interne contrôlée.
&
\textbf{Très satisfaisante}
\\
\midrule
\multicolumn{3}{r}{\textbf{Évaluation globale}} &
\textbf{Satisfaisante}
\\
\bottomrule
\end{tabularx}
\end{table}

\begin{figure}[htbp]
\centering
\begin{tikzpicture}[x=8cm,y=0.62cm]
  \draw[->,muted] (0,0) -- (0.72,0)
    node[below=4pt,font=\scriptsize] {Score moyen};
  \foreach \x/\label in {0/{0},0.2/{0,20},0.4/{0,40},0.6/{0,60}} {
    \draw[muted] (\x,0.08) -- (\x,-0.08);
    \node[below=9pt,font=\scriptsize] at (\x,0) {\label};
  }
  \draw[accent,fill=accent!75] (0,0.45) rectangle (0.153,0.80);
  \draw[accent,fill=accent!75] (0,1.45) rectangle (0.461,1.80);
  \node[anchor=east,font=\scriptsize] at (-0.05,0.625) {Français};
  \node[anchor=east,font=\scriptsize] at (-0.05,1.625) {Anglais};
  \node[anchor=west,font=\scriptsize] at (0.475,0.625) {0,153};
  \node[anchor=west,font=\scriptsize] at (0.475,1.625) {0,461};
\end{tikzpicture}
\caption{Comparaison des scores moyens selon la langue de la question}
\label{fig:langue-scores}
\end{figure}

\begin{figure}[htbp]
\centering
\begin{tikzpicture}[x=0.022cm,y=0.72cm]
  \draw[->,muted] (0,0) -- (290,0)
    node[below=4pt,font=\scriptsize] {Latence (ms)};
  \foreach \x in {0,50,100,150,200,250} {
    \draw[muted] (\x,0.08) -- (\x,-0.08);
    \node[below=9pt,font=\scriptsize] at (\x,0) {\x};
  }
  \draw[accent,fill=accent!75] (0,0.45) rectangle (193,0.80);
  \draw[accent,fill=accent!75] (0,1.45) rectangle (160,1.80);
  \draw[accent,fill=accent!75] (0,2.45) rectangle (202,2.80);
  \draw[accent,fill=accent!75] (0,3.45) rectangle (148,3.80);
  \draw[warn,fill=warn!75] (0,4.45) rectangle (263,4.80);
  \draw[accent,fill=accent!75] (0,5.45) rectangle (151,5.80);
  \draw[accent,fill=accent!75] (0,6.45) rectangle (158,6.80);
  \foreach \y/\label in {
    0.625/{Commande 10248},
    1.625/{Facture 10248},
    2.625/{Stock juillet},
    3.625/{Maria Larsson},
    4.625/{Congés payés},
    5.625/{Note de frais},
    6.625/{Intégration}
  } {
    \node[anchor=east,font=\scriptsize] at (-150,\y) {\label};
  }
  \node[anchor=west,font=\scriptsize] at (268,0.625) {193};
  \node[anchor=west,font=\scriptsize] at (235,1.625) {160};
  \node[anchor=west,font=\scriptsize] at (244,2.625) {202};
  \node[anchor=west,font=\scriptsize] at (190,3.625) {148};
  \node[anchor=west,font=\scriptsize] at (268,4.625) {263};
  \node[anchor=west,font=\scriptsize] at (193,5.625) {151};
  \node[anchor=west,font=\scriptsize] at (200,6.625) {158};
\end{tikzpicture}
\caption{Latence de recherche pour les sept requêtes évaluées}
\label{fig:latence-recherche}
\end{figure}

% ============================================================

\begin{figure}[htbp]
\centering
\begin{tikzpicture}[x=10cm,y=0.62cm]
  \draw[->,muted] (-0.1,0) -- (0.62,0)
    node[below=4pt,font=\scriptsize] {Score de pertinence maximal};
  \foreach \x/\label in {-0.1/{-0,10},0/{0},0.2/{0,20},0.4/{0,40},0.6/{0,60}} {
    \draw[muted] (\x,0.08) -- (\x,-0.08);
    \node[below=9pt,font=\scriptsize] at (\x,0) {\label};
  }
  \draw[accent,fill=accent!75] (0,0.45) rectangle (0.082,0.80);
  \draw[accent,fill=accent!75] (0,1.45) rectangle (0.115,1.80);
  \draw[accent,fill=accent!75] (0,2.45) rectangle (0.194,2.80);
  \draw[accent,fill=accent!75] (0,3.45) rectangle (0.567,3.80);
  \draw[warn,fill=warn!75] (0,4.45) rectangle (0.046,4.80);
  \draw[warn,fill=warn!75] (0,5.45) rectangle (0.003,5.80);
  \draw[warn,fill=warn!75] (-0.048,6.45) rectangle (0,6.80);
  \foreach \y/\label in {
    0.625/{Commande 10248},
    1.625/{Facture 10248},
    2.625/{Stock juillet},
    3.625/{Maria Larsson},
    4.625/{Congés payés},
    5.625/{Note de frais},
    6.625/{Intégration}
  } {
    \node[anchor=east,font=\scriptsize] at (-0.05,\y) {\label};
  }
  \node[anchor=west,font=\scriptsize] at (0.59,0.625) {0,082};
  \node[anchor=west,font=\scriptsize] at (0.59,1.625) {0,115};
  \node[anchor=west,font=\scriptsize] at (0.59,2.625) {0,194};
  \node[anchor=west,font=\scriptsize] at (0.59,3.625) {0,567};
  \node[anchor=west,font=\scriptsize] at (0.59,4.625) {0,046};
  \node[anchor=west,font=\scriptsize] at (0.59,5.625) {0,003};
  \node[anchor=west,font=\scriptsize] at (0.59,6.625) {-0,048};
\end{tikzpicture}
\caption{Scores maximaux de pertinence selon le type de question}
\label{fig:scores-retrieval}
\end{figure}

\begin{table}[htbp]
\centering
\caption{Évaluation récapitulative du système RAG}
\label{tab:evaluation-recapitulative}
\small
\begin{tabularx}{\textwidth}{@{}p{3.2cm}p{4.1cm}X p{2.1cm}@{}}
\toprule
\textbf{Dimension} & \textbf{Indicateurs} & \textbf{Résultats observés} &
\textbf{Appréciation} \\
\midrule
\shortstack[l]{\textbf{Qualité de la recherche}\\\eng{Retrieval}}
&
Score de pertinence, recherche top-$k$, couverture des sources
&
Recherche vectorielle fonctionnelle avec $k=4$ et affichage des scores.
Score moyen français : 0,153 ; score moyen anglais : 0,461. La recherche
par identifiant exact reste perfectible.
&
\textbf{Satisfaisante, à améliorer}
\\
\midrule
\shortstack[l]{\textbf{Qualité des réponses}\\\eng{Generation}}
&
Fidélité au contexte, pertinence, citations, maîtrise des hallucinations
&
Prompt restrictif, température fixée à 0,1, réponse de repli si le contexte
est absent et citation des fichiers sources. Les métriques RAGAS complètes
restent à mesurer.
&
\textbf{Satisfaisante}
\\
\midrule
\shortstack[l]{\textbf{Performance}\\\textbf{technique}}
&
Latence de recherche, latence de génération, capacité d'indexation
&
2\,676 documents traités et 2\,929 segments indexés. Latence de recherche :
182 ms en moyenne. Génération locale : environ 110 s sur CPU, ce qui
constitue le principal goulot d'étranglement.
&
\textbf{Satisfaisante}
\\
\midrule
\shortstack[l]{\textbf{Fiabilité et}\\\textbf{robustesse}}
&
Extraction, traçabilité, gestion des erreurs, reproductibilité
&
0 échec d'extraction observé sur le corpus mesuré. Métadonnées conservées,
sources et scores affichés, API structurée avec endpoints de santé,
d'interrogation et d'ingestion. Authentification et re-ranking à ajouter
avant une exposition publique.
&
\textbf{Bonne en local}
\\
\midrule
\shortstack[l]{\textbf{Utilité}\\\textbf{métier}}
&
Accès à l'information, confidentialité, simplicité d'utilisation,
vérifiabilité
&
Interface conversationnelle accessible en langage naturel, réponses
accompagnées de sources vérifiables et traitement local des documents.
L'outil est particulièrement adapté à une démonstration et à une recherche
documentaire interne contrôlée.
&
\textbf{Très satisfaisante}
\\
\midrule
\multicolumn{3}{r}{\textbf{Évaluation globale}} &
\textbf{Satisfaisante}
\\
\bottomrule
\end{tabularx}
\end{table}

\begin{figure}[htbp]
\centering
\begin{tikzpicture}[x=8cm,y=0.62cm]
  \draw[->,muted] (0,0) -- (0.72,0)
    node[below=4pt,font=\scriptsize] {Score moyen};
  \foreach \x/\label in {0/{0},0.2/{0,20},0.4/{0,40},0.6/{0,60}} {
    \draw[muted] (\x,0.08) -- (\x,-0.08);
    \node[below=9pt,font=\scriptsize] at (\x,0) {\label};
  }
  \draw[accent,fill=accent!75] (0,0.45) rectangle (0.153,0.80);
  \draw[accent,fill=accent!75] (0,1.45) rectangle (0.461,1.80);
  \node[anchor=east,font=\scriptsize] at (-0.05,0.625) {Français};
  \node[anchor=east,font=\scriptsize] at (-0.05,1.625) {Anglais};
  \node[anchor=west,font=\scriptsize] at (0.475,0.625) {0,153};
  \node[anchor=west,font=\scriptsize] at (0.475,1.625) {0,461};
\end{tikzpicture}
\caption{Comparaison des scores moyens selon la langue de la question}
\label{fig:langue-scores}
\end{figure}

\begin{figure}[htbp]
\centering
\begin{tikzpicture}[x=0.022cm,y=0.72cm]
  \draw[->,muted] (0,0) -- (290,0)
    node[below=4pt,font=\scriptsize] {Latence (ms)};
  \foreach \x in {0,50,100,150,200,250} {
    \draw[muted] (\x,0.08) -- (\x,-0.08);
    \node[below=9pt,font=\scriptsize] at (\x,0) {\x};
  }
  \draw[accent,fill=accent!75] (0,0.45) rectangle (193,0.80);
  \draw[accent,fill=accent!75] (0,1.45) rectangle (160,1.80);
  \draw[accent,fill=accent!75] (0,2.45) rectangle (202,2.80);
  \draw[accent,fill=accent!75] (0,3.45) rectangle (148,3.80);
  \draw[warn,fill=warn!75] (0,4.45) rectangle (263,4.80);
  \draw[accent,fill=accent!75] (0,5.45) rectangle (151,5.80);
  \draw[accent,fill=accent!75] (0,6.45) rectangle (158,6.80);
  \foreach \y/\label in {
    0.625/{Commande 10248},
    1.625/{Facture 10248},
    2.625/{Stock juillet},
    3.625/{Maria Larsson},
    4.625/{Congés payés},
    5.625/{Note de frais},
    6.625/{Intégration}
  } {
    \node[anchor=east,font=\scriptsize] at (-150,\y) {\label};
  }
  \node[anchor=west,font=\scriptsize] at (268,0.625) {193};
  \node[anchor=west,font=\scriptsize] at (235,1.625) {160};
  \node[anchor=west,font=\scriptsize] at (244,2.625) {202};
  \node[anchor=west,font=\scriptsize] at (190,3.625) {148};
  \node[anchor=west,font=\scriptsize] at (268,4.625) {263};
  \node[anchor=west,font=\scriptsize] at (193,5.625) {151};
  \node[anchor=west,font=\scriptsize] at (200,6.625) {158};
\end{tikzpicture}
\caption{Latence de recherche pour les sept requêtes évaluées}
\label{fig:latence-recherche}
\end{figure}

% ============================================================

\begin{figure}[htbp]
\centering
\begin{tikzpicture}[x=10cm,y=0.62cm]
  \draw[->,muted] (-0.1,0) -- (0.62,0)
    node[below=4pt,font=\scriptsize] {Score de pertinence maximal};
  \foreach \x/\label in {-0.1/{-0,10},0/{0},0.2/{0,20},0.4/{0,40},0.6/{0,60}} {
    \draw[muted] (\x,0.08) -- (\x,-0.08);
    \node[below=9pt,font=\scriptsize] at (\x,0) {\label};
  }
  \draw[accent,fill=accent!75] (0,0.45) rectangle (0.082,0.80);
  \draw[accent,fill=accent!75] (0,1.45) rectangle (0.115,1.80);
  \draw[accent,fill=accent!75] (0,2.45) rectangle (0.194,2.80);
  \draw[accent,fill=accent!75] (0,3.45) rectangle (0.567,3.80);
  \draw[warn,fill=warn!75] (0,4.45) rectangle (0.046,4.80);
  \draw[warn,fill=warn!75] (0,5.45) rectangle (0.003,5.80);
  \draw[warn,fill=warn!75] (-0.048,6.45) rectangle (0,6.80);
  \foreach \y/\label in {
    0.625/{Commande 10248},
    1.625/{Facture 10248},
    2.625/{Stock juillet},
    3.625/{Maria Larsson},
    4.625/{Congés payés},
    5.625/{Note de frais},
    6.625/{Intégration}
  } {
    \node[anchor=east,font=\scriptsize] at (-0.05,\y) {\label};
  }
  \node[anchor=west,font=\scriptsize] at (0.59,0.625) {0,082};
  \node[anchor=west,font=\scriptsize] at (0.59,1.625) {0,115};
  \node[anchor=west,font=\scriptsize] at (0.59,2.625) {0,194};
  \node[anchor=west,font=\scriptsize] at (0.59,3.625) {0,567};
  \node[anchor=west,font=\scriptsize] at (0.59,4.625) {0,046};
  \node[anchor=west,font=\scriptsize] at (0.59,5.625) {0,003};
  \node[anchor=west,font=\scriptsize] at (0.59,6.625) {-0,048};
\end{tikzpicture}
\caption{Scores maximaux de pertinence selon le type de question}
\label{fig:scores-retrieval}
\end{figure}

\begin{table}[htbp]
\centering
\caption{Évaluation récapitulative du système RAG}
\label{tab:evaluation-recapitulative}
\small
\begin{tabularx}{\textwidth}{@{}p{3.2cm}p{4.1cm}X p{2.1cm}@{}}
\toprule
\textbf{Dimension} & \textbf{Indicateurs} & \textbf{Résultats observés} &
\textbf{Appréciation} \\
\midrule
\shortstack[l]{\textbf{Qualité de la recherche}\\\eng{Retrieval}}
&
Score de pertinence, recherche top-$k$, couverture des sources
&
Recherche vectorielle fonctionnelle avec $k=4$ et affichage des scores.
Score moyen français : 0,153 ; score moyen anglais : 0,461. La recherche
par identifiant exact reste perfectible.
&
\textbf{Satisfaisante, à améliorer}
\\
\midrule
\shortstack[l]{\textbf{Qualité des réponses}\\\eng{Generation}}
&
Fidélité au contexte, pertinence, citations, maîtrise des hallucinations
&
Prompt restrictif, température fixée à 0,1, réponse de repli si le contexte
est absent et citation des fichiers sources. Les métriques RAGAS complètes
restent à mesurer.
&
\textbf{Satisfaisante}
\\
\midrule
\shortstack[l]{\textbf{Performance}\\\textbf{technique}}
&
Latence de recherche, latence de génération, capacité d'indexation
&
2\,676 documents traités et 2\,929 segments indexés. Latence de recherche :
182 ms en moyenne. Génération locale : environ 110 s sur CPU, ce qui
constitue le principal goulot d'étranglement.
&
\textbf{Satisfaisante}
\\
\midrule
\shortstack[l]{\textbf{Fiabilité et}\\\textbf{robustesse}}
&
Extraction, traçabilité, gestion des erreurs, reproductibilité
&
0 échec d'extraction observé sur le corpus mesuré. Métadonnées conservées,
sources et scores affichés, API structurée avec endpoints de santé,
d'interrogation et d'ingestion. Authentification et re-ranking à ajouter
avant une exposition publique.
&
\textbf{Bonne en local}
\\
\midrule
\shortstack[l]{\textbf{Utilité}\\\textbf{métier}}
&
Accès à l'information, confidentialité, simplicité d'utilisation,
vérifiabilité
&
Interface conversationnelle accessible en langage naturel, réponses
accompagnées de sources vérifiables et traitement local des documents.
L'outil est particulièrement adapté à une démonstration et à une recherche
documentaire interne contrôlée.
&
\textbf{Très satisfaisante}
\\
\midrule
\multicolumn{3}{r}{\textbf{Évaluation globale}} &
\textbf{Satisfaisante}
\\
\bottomrule
\end{tabularx}
\end{table}

\begin{figure}[htbp]
\centering
\begin{tikzpicture}[x=8cm,y=0.62cm]
  \draw[->,muted] (0,0) -- (0.72,0)
    node[below=4pt,font=\scriptsize] {Score moyen};
  \foreach \x/\label in {0/{0},0.2/{0,20},0.4/{0,40},0.6/{0,60}} {
    \draw[muted] (\x,0.08) -- (\x,-0.08);
    \node[below=9pt,font=\scriptsize] at (\x,0) {\label};
  }
  \draw[accent,fill=accent!75] (0,0.45) rectangle (0.153,0.80);
  \draw[accent,fill=accent!75] (0,1.45) rectangle (0.461,1.80);
  \node[anchor=east,font=\scriptsize] at (-0.05,0.625) {Français};
  \node[anchor=east,font=\scriptsize] at (-0.05,1.625) {Anglais};
  \node[anchor=west,font=\scriptsize] at (0.475,0.625) {0,153};
  \node[anchor=west,font=\scriptsize] at (0.475,1.625) {0,461};
\end{tikzpicture}
\caption{Comparaison des scores moyens selon la langue de la question}
\label{fig:langue-scores}
\end{figure}

\begin{figure}[htbp]
\centering
\begin{tikzpicture}[x=0.022cm,y=0.72cm]
  \draw[->,muted] (0,0) -- (290,0)
    node[below=4pt,font=\scriptsize] {Latence (ms)};
  \foreach \x in {0,50,100,150,200,250} {
    \draw[muted] (\x,0.08) -- (\x,-0.08);
    \node[below=9pt,font=\scriptsize] at (\x,0) {\x};
  }
  \draw[accent,fill=accent!75] (0,0.45) rectangle (193,0.80);
  \draw[accent,fill=accent!75] (0,1.45) rectangle (160,1.80);
  \draw[accent,fill=accent!75] (0,2.45) rectangle (202,2.80);
  \draw[accent,fill=accent!75] (0,3.45) rectangle (148,3.80);
  \draw[warn,fill=warn!75] (0,4.45) rectangle (263,4.80);
  \draw[accent,fill=accent!75] (0,5.45) rectangle (151,5.80);
  \draw[accent,fill=accent!75] (0,6.45) rectangle (158,6.80);
  \foreach \y/\label in {
    0.625/{Commande 10248},
    1.625/{Facture 10248},
    2.625/{Stock juillet},
    3.625/{Maria Larsson},
    4.625/{Congés payés},
    5.625/{Note de frais},
    6.625/{Intégration}
  } {
    \node[anchor=east,font=\scriptsize] at (-150,\y) {\label};
  }
  \node[anchor=west,font=\scriptsize] at (268,0.625) {193};
  \node[anchor=west,font=\scriptsize] at (235,1.625) {160};
  \node[anchor=west,font=\scriptsize] at (244,2.625) {202};
  \node[anchor=west,font=\scriptsize] at (190,3.625) {148};
  \node[anchor=west,font=\scriptsize] at (268,4.625) {263};
  \node[anchor=west,font=\scriptsize] at (193,5.625) {151};
  \node[anchor=west,font=\scriptsize] at (200,6.625) {158};
\end{tikzpicture}
\caption{Latence de recherche pour les sept requêtes évaluées}
\label{fig:latence-recherche}
\end{figure}

% ============================================================

\begin{figure}[htbp]
\centering
\begin{tikzpicture}[x=10cm,y=0.62cm]
  \draw[->,muted] (-0.1,0) -- (0.62,0)
    node[below=4pt,font=\scriptsize] {Score de pertinence maximal};
  \foreach \x/\label in {-0.1/{-0,10},0/{0},0.2/{0,20},0.4/{0,40},0.6/{0,60}} {
    \draw[muted] (\x,0.08) -- (\x,-0.08);
    \node[below=9pt,font=\scriptsize] at (\x,0) {\label};
  }
  \draw[accent,fill=accent!75] (0,0.45) rectangle (0.082,0.80);
  \draw[accent,fill=accent!75] (0,1.45) rectangle (0.115,1.80);
  \draw[accent,fill=accent!75] (0,2.45) rectangle (0.194,2.80);
  \draw[accent,fill=accent!75] (0,3.45) rectangle (0.567,3.80);
  \draw[warn,fill=warn!75] (0,4.45) rectangle (0.046,4.80);
  \draw[warn,fill=warn!75] (0,5.45) rectangle (0.003,5.80);
  \draw[warn,fill=warn!75] (-0.048,6.45) rectangle (0,6.80);
  \foreach \y/\label in {
    0.625/{Commande 10248},
    1.625/{Facture 10248},
    2.625/{Stock juillet},
    3.625/{Maria Larsson},
    4.625/{Congés payés},
    5.625/{Note de frais},
    6.625/{Intégration}
  } {
    \node[anchor=east,font=\scriptsize] at (-0.05,\y) {\label};
  }
  \node[anchor=west,font=\scriptsize] at (0.59,0.625) {0,082};
  \node[anchor=west,font=\scriptsize] at (0.59,1.625) {0,115};
  \node[anchor=west,font=\scriptsize] at (0.59,2.625) {0,194};
  \node[anchor=west,font=\scriptsize] at (0.59,3.625) {0,567};
  \node[anchor=west,font=\scriptsize] at (0.59,4.625) {0,046};
  \node[anchor=west,font=\scriptsize] at (0.59,5.625) {0,003};
  \node[anchor=west,font=\scriptsize] at (0.59,6.625) {-0,048};
\end{tikzpicture}
\caption{Scores maximaux de pertinence selon le type de question}
\label{fig:scores-retrieval}
\end{figure}

\begin{table}[htbp]
\centering
\caption{Évaluation récapitulative du système RAG}
\label{tab:evaluation-recapitulative}
\small
\begin{tabularx}{\textwidth}{@{}p{3.2cm}p{4.1cm}X p{2.1cm}@{}}
\toprule
\textbf{Dimension} & \textbf{Indicateurs} & \textbf{Résultats observés} &
\textbf{Appréciation} \\
\midrule
\shortstack[l]{\textbf{Qualité de la recherche}\\\eng{Retrieval}}
&
Score de pertinence, recherche top-$k$, couverture des sources
&
Recherche vectorielle fonctionnelle avec $k=4$ et affichage des scores.
Score moyen français : 0,153 ; score moyen anglais : 0,461. La recherche
par identifiant exact reste perfectible.
&
\textbf{Satisfaisante, à améliorer}
\\
\midrule
\shortstack[l]{\textbf{Qualité des réponses}\\\eng{Generation}}
&
Fidélité au contexte, pertinence, citations, maîtrise des hallucinations
&
Prompt restrictif, température fixée à 0,1, réponse de repli si le contexte
est absent et citation des fichiers sources. Les métriques RAGAS complètes
restent à mesurer.
&
\textbf{Satisfaisante}
\\
\midrule
\shortstack[l]{\textbf{Performance}\\\textbf{technique}}
&
Latence de recherche, latence de génération, capacité d'indexation
&
2\,676 documents traités et 2\,929 segments indexés. Latence de recherche :
182 ms en moyenne. Génération locale : environ 110 s sur CPU, ce qui
constitue le principal goulot d'étranglement.
&
\textbf{Satisfaisante}
\\
\midrule
\shortstack[l]{\textbf{Fiabilité et}\\\textbf{robustesse}}
&
Extraction, traçabilité, gestion des erreurs, reproductibilité
&
0 échec d'extraction observé sur le corpus mesuré. Métadonnées conservées,
sources et scores affichés, API structurée avec endpoints de santé,
d'interrogation et d'ingestion. Authentification et re-ranking à ajouter
avant une exposition publique.
&
\textbf{Bonne en local}
\\
\midrule
\shortstack[l]{\textbf{Utilité}\\\textbf{métier}}
&
Accès à l'information, confidentialité, simplicité d'utilisation,
vérifiabilité
&
Interface conversationnelle accessible en langage naturel, réponses
accompagnées de sources vérifiables et traitement local des documents.
L'outil est particulièrement adapté à une démonstration et à une recherche
documentaire interne contrôlée.
&
\textbf{Très satisfaisante}
\\
\midrule
\multicolumn{3}{r}{\textbf{Évaluation globale}} &
\textbf{Satisfaisante}
\\
\bottomrule
\end{tabularx}
\end{table}

\begin{figure}[htbp]
\centering
\begin{tikzpicture}[x=8cm,y=0.62cm]
  \draw[->,muted] (0,0) -- (0.72,0)
    node[below=4pt,font=\scriptsize] {Score moyen};
  \foreach \x/\label in {0/{0},0.2/{0,20},0.4/{0,40},0.6/{0,60}} {
    \draw[muted] (\x,0.08) -- (\x,-0.08);
    \node[below=9pt,font=\scriptsize] at (\x,0) {\label};
  }
  \draw[accent,fill=accent!75] (0,0.45) rectangle (0.153,0.80);
  \draw[accent,fill=accent!75] (0,1.45) rectangle (0.461,1.80);
  \node[anchor=east,font=\scriptsize] at (-0.05,0.625) {Français};
  \node[anchor=east,font=\scriptsize] at (-0.05,1.625) {Anglais};
  \node[anchor=west,font=\scriptsize] at (0.475,0.625) {0,153};
  \node[anchor=west,font=\scriptsize] at (0.475,1.625) {0,461};
\end{tikzpicture}
\caption{Comparaison des scores moyens selon la langue de la question}
\label{fig:langue-scores}
\end{figure}

\begin{figure}[htbp]
\centering
\begin{tikzpicture}[x=0.022cm,y=0.72cm]
  \draw[->,muted] (0,0) -- (290,0)
    node[below=4pt,font=\scriptsize] {Latence (ms)};
  \foreach \x in {0,50,100,150,200,250} {
    \draw[muted] (\x,0.08) -- (\x,-0.08);
    \node[below=9pt,font=\scriptsize] at (\x,0) {\x};
  }
  \draw[accent,fill=accent!75] (0,0.45) rectangle (193,0.80);
  \draw[accent,fill=accent!75] (0,1.45) rectangle (160,1.80);
  \draw[accent,fill=accent!75] (0,2.45) rectangle (202,2.80);
  \draw[accent,fill=accent!75] (0,3.45) rectangle (148,3.80);
  \draw[warn,fill=warn!75] (0,4.45) rectangle (263,4.80);
  \draw[accent,fill=accent!75] (0,5.45) rectangle (151,5.80);
  \draw[accent,fill=accent!75] (0,6.45) rectangle (158,6.80);
  \foreach \y/\label in {
    0.625/{Commande 10248},
    1.625/{Facture 10248},
    2.625/{Stock juillet},
    3.625/{Maria Larsson},
    4.625/{Congés payés},
    5.625/{Note de frais},
    6.625/{Intégration}
  } {
    \node[anchor=east,font=\scriptsize] at (-150,\y) {\label};
  }
  \node[anchor=west,font=\scriptsize] at (268,0.625) {193};
  \node[anchor=west,font=\scriptsize] at (235,1.625) {160};
  \node[anchor=west,font=\scriptsize] at (244,2.625) {202};
  \node[anchor=west,font=\scriptsize] at (190,3.625) {148};
  \node[anchor=west,font=\scriptsize] at (268,4.625) {263};
  \node[anchor=west,font=\scriptsize] at (193,5.625) {151};
  \node[anchor=west,font=\scriptsize] at (200,6.625) {158};
\end{tikzpicture}
\caption{Latence de recherche pour les sept requêtes évaluées}
\label{fig:latence-recherche}
\end{figure}
